%% ma: L12-B08-0123
%% lop: 12
%% chuong: 2
%% bai: 8
%% dang: 12.8.45
%% loai: TN4
%% muc: VD
%% dap_an: "A"
%% nguon: "CĐ2.3.pdf giáo viên gửi 27/10/2026 (M07, Lê Minh Khôi – Nguyễn Anh Khôi, Chương II Chủ đề 3 Biểu thức tọa độ của các phép toán vectơ), Câu 36"
%% kien_thuc: "Bài 8 (tích vô hướng, khoảng cách hai điểm), lớp 9 (tâm đường tròn ngoại tiếp tam giác vuông)"
%% trang_thai: chinh-thuc
%% ghi_chu: "Giáo viên đã duyệt 27/10/2026"
\begin{ex}
Trong không gian $Oxyz$, cho hai điểm $A(1;-2;3)$, $B(-2;2;2)$. Gọi $I(a;b;c)$ là tâm đường tròn ngoại tiếp tam giác $OAB$. Tính giá trị của biểu thức $T=a^2+b^2+c^2$.
\choice{\True $T=\dfrac{13}2$}{$T=2$}{$T=6$}{$T=\dfrac{29}4$}
\loigiai{$\overrightarrow{OA}\cdot\overrightarrow{OB}=-2-4+6=0$ nên tam giác $OAB$ vuông tại $O$, tâm đường tròn ngoại tiếp là trung điểm $AB$.

$I\left(-\dfrac12;0;\dfrac52\right)$, $T=\dfrac14+\dfrac{25}4=\dfrac{13}2$.}
\end{ex}
